%% titre: Méthodes — Étudier la convexité
%% matiere: Optimisation
%% niveau: Licence 3
%% annee: 2026-2027
%% ordre: 2
%% description: Convexité par la hessienne, critère de Sylvester, fonctionnelles quadratiques.

\section{Convexité d'une fonction de classe $\mathcal C^2$}

\begin{methode}[Par la hessienne]
Soit $f$ de classe $\mathcal C^2$ sur un ouvert convexe $C$.
\begin{enumerate}[label=\textbf{Étape \arabic*.}]
  \item Calculer $\nabla^2f(x)$ pour tout $x\in C$.
  \item Étudier son signe : $\nabla^2f(x)\succeq0$ pour tout $x$ $\iff$ $f$ convexe.
  \item Si $\nabla^2f(x)\succ0$ pour tout $x$, alors $f$ est strictement convexe.
\end{enumerate}
\end{methode}

\begin{attention}
La réciproque de la dernière implication est fausse : $x\mapsto x^4$ est strictement convexe mais $f''(0)=0$.
\end{attention}

\section{Signe d'une matrice : Sylvester}

\begin{methode}[Critère de Sylvester]
Pour $A$ symétrique, calculer les mineurs principaux dominants $\Delta_1,\dots,\Delta_n$ (déterminants des sous-matrices du coin supérieur gauche).
\begin{itemize}
  \item $A\succ0$ $\iff$ $\Delta_k>0$ pour tout $k$.
  \item $A\prec0$ $\iff$ $(-1)^k\Delta_k>0$ pour tout $k$.
  \item Si certains $\Delta_k$ sont nuls : on ne peut rien conclure pour le caractère \emph{semi}-défini. Calculer les valeurs propres ($A\succeq0\iff$ toutes $\ge0$).
\end{itemize}
\end{methode}

\begin{exemple}
Pour $f(x,y)=x^2+xy+y^2$, $\nabla^2f=\begin{pmatrix}2&1\\1&2\end{pmatrix}$ : $\Delta_1=2>0$ et $\Delta_2=3>0$. La hessienne est définie positive : $f$ est strictement convexe.

Pour $g(x,y)=x^2+4xy+y^2$, $\nabla^2g=\begin{pmatrix}2&4\\4&2\end{pmatrix}$ : $\Delta_2=4-16=-12<0$. La hessienne est indéfinie : $g$ n'est pas convexe.
\end{exemple}

\section{Fonctionnelle quadratique}

\begin{methode}[$f(x)=\tfrac12x^{\top}Ax-b^{\top}x+c$]
\begin{enumerate}[label=\textbf{Étape \arabic*.}]
  \item Symétriser si besoin : remplacer $A$ par $\frac12(A+A^{\top})$.
  \item Écrire $\nabla f(x)=Ax-b$ et $\nabla^2f(x)=A$ (hessienne \emph{constante}).
  \item Conclure : $f$ convexe $\iff A\succeq0$ ; strictement convexe $\iff A\succ0$.
  \item Si $A\succ0$ : l'unique minimum global est $x^*=A^{-1}b$.
\end{enumerate}
\end{methode}

\begin{exemple}
Soit $f(x,y,z)=x^2+2y^2+3z^2+2xy+2yz=\tfrac12v^{\top}Av$ avec $A=\begin{pmatrix}2&2&0\\2&4&2\\0&2&6\end{pmatrix}$. On a $\Delta_1=2$, $\Delta_2=8-4=4$, $\Delta_3=16$, tous $>0$ : $A\succ0$, $f$ est strictement convexe et son unique minimum global est $(0,0,0)$.
\end{exemple}

\begin{exercice}[Application]
Étudier la convexité de $f(x,y)=x^2-2xy+y^2$.
\begin{correction}
$\nabla^2f=\begin{pmatrix}2&-2\\-2&2\end{pmatrix}$ avec $\Delta_1=2>0$ mais $\Delta_2=0$ : Sylvester ne permet pas de conclure. Les valeurs propres valent $0$ et $4$, toutes $\ge0$ : $\nabla^2f\succeq0$ et $f$ est convexe, mais pas strictement (on a $f(x,y)=(x-y)^2$, nulle sur toute la droite $y=x$).
\end{correction}
\end{exercice}
